\question{Which model and protocol should the pilot deploy, and why does that matter legally?}
\label{sec:modelchoice}
\statusMeasured

\begin{shortanswer}
Logistic regression under FedProx or FedAvg. The legally useful fact is that this recommendation is
\textbf{over-determined}: the same choice wins on accuracy, on bandwidth, on disclosure surface, and
on protectability. There is no trade-off here in which privacy is bought with performance the
model that federates best is also the only one secure aggregation can protect.
\end{shortanswer}

\subsection{Three findings, one conclusion}

\paragraph{FedProx and FedAvg are indistinguishable.} They agree to three decimals on both measures:
\BestFedAUC{} against \FedAvgAUC{} global, and \BestFedWorst{} worst practice for both. FedProx adds a proximal term
\cite{li2020fedprox} intended to correct client drift, and at two local passes per round there is
almost no drift to correct. This is a null result and is reported as one; it should be retested if
the local epoch count increases, since that is the condition under which the term would start to
matter.

\paragraph{FedMosaic does not win on accuracy, and its case is elsewhere.} It reaches \MosaicAUC{}
global slightly below the weight-sharing protocols, at \UplinkMosaic{} bits per round against
\UplinkLogreg{}. Its argument is its disclosure profile: it
never transmits a model at all, only predictions on an already-public cohort (\S\ref{sec:payload}).
That is a qualitative difference, and if the legal analysis ends up treating parameter vectors as
personal data while treating predictions on public patients differently, FedMosaic becomes worth
revisiting on those grounds not on these.

\paragraph{The tree-based alternative was evaluated and is no longer in scope.} A
gradient-boosted-trees arm (FedXgbBagging) was benchmarked against the logistic protocols on the
pinned run this report is built from: its pooled ceiling matched logistic regression, so the model
class was not the problem, but federated it reached far less, it \emph{degraded} with additional
rounds (\S\ref{sec:reduce}), its split thresholds were literal patient values
(\S\ref{sec:payload}), and secure aggregation could not protect it at all (\S\ref{sec:secagg}).
The measured comparison is retained, marked historical, in the accompanying benchmark report
(the test bed's \texttt{docs/REPORT.md}, scope note 2026-08-28). On 2026-08-28 the tree-based and MLP
code paths were \textbf{deleted from the test bed}; logistic regression --- the model named in
grant task T2.2 --- is now the only model class in the codebase, so the recommendation is not a
preference among alternatives but a description of what exists.

\subsection{Why the convergence matters to counsel}

It is common for privacy measures to be resisted on the grounds that they cost accuracy, which turns
every safeguard into a negotiation. That is not the situation here. The evidence points one way on
every axis at once (numbers as measured on the pinned run; the ``against'' column is the
now-removed trees arm, retained as historical context):

\begin{center}
\small
\begin{tabular}{lll}
  \toprule
  Criterion & Favours & Against (historical) \\
  \midrule
  Global accuracy & logistic regression & trees (removed 2026-08-28) \\
  Worst-practice accuracy & logistic regression & trees (removed 2026-08-28) \\
  Behaviour with more rounds & logistic regression (flat) & trees (degrade) \\
  Bandwidth & logistic regression & trees ($\sim$60$\times$ the uplink) \\
  Disclosure surface & logistic regression & trees (literal patient values) \\
  Protectable by secure aggregation & logistic regression & trees (impossible) \\
  \bottomrule
\end{tabular}
\end{center}

The consortium therefore does not have to trade model performance for the ability to make privacy
guarantees. Pinning the deployment to logistic regression is what makes secure aggregation possible
at all, and it costs nothing to do so.
